% chapters/11/checklist.tex
\bagian{Checklist Pencapaian Kompetensi}

Bab ini menuntaskan \SUBCPMK{2.4} dan menutup \CPMK{2}. Centang setiap butir
berikut hanya apabila Anda benar-benar telah melakukannya atau dapat
menjelaskannya.

\begin{checklistbox}
\begin{checklist}
  \item Saya dapat menjelaskan perbedaan pekerjaan lapisan keputusan dan
        lapisan navigasi, serta cara keduanya dihubungkan.
  \item Saya dapat menyimpan hasil pencarian jalur dan memakainya berkali-kali.
  \item Saya dapat menjelaskan dua sebab pencarian jalur perlu diulang.
  \item Saya dapat menulis pengikut jalur yang memakai perulangan
        \code{while} dan toleransi jarak, beserta alasan keduanya.
  \item {\bfseries (Indikator a)} Saya dapat menyusun skenario uji yang dapat
        diulang dan memicu perilaku yang hendak dinilai.
  \item Saya telah menguji mini proyek pada ketiga skenario yang diminta,
        yaitu ruang terbuka, labirin, dan rintangan dinamis.
  \item Laporan \istilah{playtest} saya memuat angka pengukuran dan kesan mutu
        perilaku pada setiap skenario.
  \item {\bfseries (Indikator b)} Saya dapat menjelaskan tiga segi penilaian
        mutu perilaku NPC beserta ukuran masing-masing.
  \item Saya telah menjalankan evaluasi selang hitung ulang dan mencatat
        hasilnya pada lembar kerja.
  \item Saya dapat menghitung rasio beban dan selisih mutu antara dua selang
        hitung ulang, lalu menyimpulkan mana yang lebih menguntungkan.
  \item Saya dapat menjelaskan mengapa menghitung ulang setiap bingkai
        merupakan pemborosan.
  \item Saya dapat menjelaskan mengapa selang yang terlalu panjang merugikan
        perilaku NPC.
  \item {\bfseries (Indikator c)} Saya telah menerapkan sekurang-kurangnya satu
        perbaikan dari tabel perbaikan, dan mengukur hasilnya.
  \item Saya dapat menjelaskan mengapa pekerjaan yang mahal tetapi tidak
        mendesak sebaiknya disebar, bukan ditumpuk.
  \item Mini proyek kelompok saya menghasilkan kode, ringkasan desain, laporan
        \istilah{playtest}, dan presentasi.
  \item Setiap anggota kelompok saya memahami bagiannya masing-masing dan
        dapat menjawab pertanyaan penguji.
\end{checklist}
\end{checklistbox}
